\subsection{Toroidal incidence, residual cocycle, and fixture validation}
\label{sec:bt1514-bt1516-toroid-cocycle-fixture}

The 7/21/3 bridge passes a Szilassi-style incidence-count test: seven face classes, twenty-one edge/flag classes, six incident edges per face, two incident faces per edge, and three two-edge local sectors per face.  This is incidence compatibility, not yet a concrete realization theorem.

The residual key now has an algebraic gauge-cocycle law.  Under left and right line-gauge changes the residual transforms as \(\rho\mapsto R^{-1}\rho L\).  The finite \(S_3\) key action is closed, invertible, and associative.

The native D4 calibration fixture schema validates against route-trace and CSS-replay shapes: eight native generators, seventy-two rows per generator, 576 generated rows, 192 active rows, and 384 guard rows.
